\section{Introduction}

A formal proof often changes the language in which its objects are
represented. A graph becomes a Boolean matrix; an algebraic wedge product
becomes a continuous bilinear operator; a randomized estimator becomes a
program reading addressed tape bits. The next argument needs more than a
new description of the object. It needs a theorem preserving the
particular quantity or operation on which that argument depends.

We call such correspondences preservation interfaces. An interface
specifies the relation between two representations and exposes the
information available to a later construction. To transfer an accuracy
guarantee, it must preserve the relevant event probability; equality of
output laws supplies a reusable sufficient condition. To differentiate a
field, it must provide appropriate local agreement. To transport an
equation in named coordinates, it must identify the coordinate maps.
These requirements explain both the mathematical content and the formal
shape of the interfaces studied here.

Six developments in the OpenAI Lean collection provide the comparison
\citep{OAIRepository}. Common-base counting and graph switch chains
connect finite representations to probabilistic guarantees. The Mahler
and Crouzeix arguments connect algebraic identities to differentiation,
integration, and norm estimates. Trace ideals and marked character-variety
diagrams reconstruct or transport algebraic data. The first four selected
statements are public theorem endpoints; the last two are support results.
Keeping these roles explicit allows the comparison to follow the
mathematical work performed by each selected construction.

Three questions guide each case: what does the intermediate representation
retain, which theorem relates it to the intended object, and how is that
theorem used in the next step? The case studies answer these questions
through source declarations and worked mathematical examples. A separate
exploratory survey compares the written proof methods of 800 uniformly
sampled OpenAI modules with the pinned Mathlib dependency. It distinguishes
average within-module practice from body frequencies dominated by large
certificate batches, then examines broad subject and proof-length
composition. The two analyses serve mathematicians reading formal source
and formalizers designing reusable interfaces.

Together, they describe a division of proof work. Representation and
preservation lemmas make the hypotheses of a local calculation available;
simplification, arithmetic tactics, or exact decision procedures then
check the resulting obligation. Arithmetic and algebra tactics remain
substantially more prevalent after joint subject and length controls.
The corresponding differences in tactic blocks and intermediate claims
are smaller under these controls. Source examples explain how such
devices operate within particular proofs: the construction determines
which goals reach a tactic and why the available hypotheses suffice.

Companion reviews assess the proposed mathematical advances against
prior results \citep{GounderFrontiers2026} and examine counting and
sampling applications \citep{GounderCounting2026}. The present review
uses common bases, Mahler, and Crouzeix to study a different question:
how the formal source connects its mathematical representations and
organizes the proof obligations between them.

The comparison builds on established accounts of formalization. Lean
combines dependent type theory with an extensible elaborator and a
programming language \citep{deMouraUllrich2021}; Mathlib supplies a shared
mathematical vocabulary \citep{mathlib2020}.
\citet{AvigadEtAl2024Anatomy} trace a linear-algebra example through
formal proof, generalization, and connection to an existing Mathlib
theorem. \citet{CommelinTopaz2024} explain how usable specifications
support mathematical construction and collaboration. Here we compare
the mathematical payload of preservation statements across several
settings and relate their local uses to a corpus-level method profile.

The odd-order theorem and Kepler conjecture establish precedents for
large proof decompositions and verified computation
\citep{GonthierEtAl2013,HalesEtAl2017}. Lean formalizations of perfectoid
spaces and schemes illustrate the importance of representation choices
\citep{Perfectoid2020,Schemes2021}. These precedents frame the novelty
question in this review: which established device is used, what
project-specific information does its interface preserve, and what later
step becomes possible because that information is available?

Section~\ref{sec:method} describes selection and evidence.
Section~\ref{sec:cases} begins with the overview in
Table~\ref{tab:overview}, then develops the six cases.
Section~\ref{sec:profile} compares source-written proof methods with
Mathlib, including joint composition controls and certificate
concentration. Section~\ref{sec:worked} isolates the cases' preservation
requirements in elementary examples. Section~\ref{sec:discussion}
connects the mathematical and empirical findings and develops guidance
for new formalizations.
